\section*{Abstract}
\begin{abstract}
  \noindent
  Dieser Projektbericht untersucht Cross-Site Scripting (XSS) und
  Cross-Origin Resource Sharing (CORS) anhand der realen
  Vue-/Spring-Boot-Anwendung \emph{Scheduler}. Die Themen wurden
  innerhalb der Gruppe aufgeteilt: Arvid Freese und Linus Paliga
  bearbeiteten XSS, Tom Markmann und Raik Ole Kannenberg CORS. Nach
  der Herleitung der theoretischen Grundlagen wurde das
  Scheduler-Repository systematisch auf relevante Datenflüsse,
  Ausgabesenken und Sicherheitskonfigurationen geprüft. Für beide
  Themen entstanden anschließend lokale Live-Demos.

  Die XSS-Analyse identifizierte eine reale
  Stored-DOM-XSS-Schwachstelle. Datenbankbasierte
  Veranstaltungsnamen, Gruppenbezeichnungen und Reservierungslabels
  wurden in mehreren Kalenderansichten an den Standardrenderer der
  Bibliothek \texttt{vue-cal} übergeben, der Event-Titel und -Inhalte
  als HTML rendert. Die Schwachstelle wurde durch eigene Event-Slots
  mit Vue-Interpolation strukturell behoben. Zum Berichtsstand ist
  dieser Fix auf \texttt{capstone-main} in allen relevanten
  Kalenderansichten ausgerollt; nach erneuter Prüfung besteht kein
  bekanntes aktuelles XSS-Risiko dieser Art.

  Die Spring-API besaß im untersuchten Ausgangsstand keine pauschale
  CORS-Freigabe. Die normale lokale Anwendung vermeidet
  Cross-Origin-Aufrufe über den Vite-Proxy. Für die Demo wurde im
  separaten Security-Branch bewusst eine unsichere globale
  Origin-Freigabe ergänzt und anschließend einer festen Allowlist
  gegenübergestellt. Unabhängig davon bleiben konkrete Verbesserungen
  erforderlich: \texttt{/files/**} ist zu breit öffentlich
  freigegeben, Keycloak verwendet für den Web-Client eine
  Wildcard-Origin sowie nicht benötigte OAuth-Flows, und das Frontend
  setzt trotz Bearer-Token-Authentifizierung \texttt{credentials:
  include}. Der Bericht leitet daraus priorisierte, direkt auf
  Scheduler-Dateien bezogene Maßnahmen ab.
\end{abstract}
